\documentclass[11pt]{article}
\usepackage[utf8]{inputenc}
\usepackage{amsthm,amsmath,amsfonts}
\input wims.sty
\theoremstyle{plain}
\newtheorem{thm}{Théorème}[section]
\newtheorem{prop}[thm]{Proposition}
\theoremstyle{definition}
\newtheorem{defn}{Définition}[section]
\newtheorem{exo}{Exercice}
\newtheorem{rem}{Remarque}
\newenvironment{sol}{}{}
\begin{document}
\title{Environnements wims.sty}
\author{Corpus latex2wims}
\email{corpus@exemple.fr}
\section{Théorèmes et démonstrations}
\begin{thm}[Pythagore]
Dans un triangle rectangle, \(a^2+b^2=c^2\).
\end{thm}
\begin{proof}
Découper le carré construit sur l'hypoténuse.
\end{proof}
\begin{defn}
Une suite \((u_n)\) est \emph{croissante} si \(u_{n+1}\geq u_n\).
\end{defn}
\begin{rem}
Remarque avec \textbf{gras}, \textit{italique} et \texttt{code}.
\end{rem}
\subsection{Plis (typefold)}
\begin{exo}
Calculer \(\sum_{k=1}^n k\).
\begin{sol}
\[ \sum_{k=1}^n k = \frac{n(n+1)}{2}. \]
\end{sol}
\end{exo}
\section{Listes et tableaux}
\begin{itemize}
\item premier élément ;
\item second élément, avec \(\sqrt{2}\).
\end{itemize}
\begin{enumerate}
\item un
\item deux
\end{enumerate}
\begin{description}
\item[Graphe] un ensemble de sommets et d'arêtes.
\item[Arbre] un graphe connexe sans cycle.
\end{description}
\begin{tabular}{|l|c|}
\hline
\(x\) & \(x^2\) \\
\hline
2 & 4 \\
\hline
\end{tabular}
\section{Liens WIMS}
\begin{center}
Un exercice : \exercise{module=H6/analysis/oefcourbe.fr&exo=courb1}{Courbes}.
\end{center}
\begin{quote}
Une citation.
\end{quote}
Un lien externe : \url{https://wims.univ-cotedazur.fr}.
\begin{latexonly}
Ce paragraphe n'apparaît que dans le PDF.
\end{latexonly}
\begin{wimsonly}
Ce paragraphe n'apparaît que dans WIMS.
\end{wimsonly}
\end{document}
